\section{The Target Open-Source Ecosystem}
\label{sec:ecosystem}

Track~3 asks proposers to improve the safety, security, and privacy of an
\emph{existing} open-source ecosystem. This project targets the ecosystem
that forms the observable and defensible perimeter of nearly every
open-source deployment, anchored on two widely used, publicly available
products: \textbf{nginx}~\cite{nginx,nginxsrc}, through which traffic enters
and moves within such deployments, and \textbf{ZMap}~\cite{zmap,zmapsrc},
with which operators discover what is actually running. \sysname hardens the
first, extends the second, and uses both as the concrete ecosystem against
which its methods are developed and evaluated.

\paragraph{nginx: the trust boundary of open-source infrastructure.}
nginx is an HTTP server, reverse proxy, load balancer, and TLS terminator
that has been released under the two-clause BSD license and that has been in
continuous development since 2004~\cite{nginx}. It ranks consistently among
the most widely deployed web servers in the world. The
\emph{development and testing model} follows a maintainer-review process
around a public repository~\cite{nginxsrc} and a development mailing list,
with periodic mainline and stable branches, an in-tree functional test suite,
continuous integration across a broad matrix of platforms and build
configurations, and downstream testing by the distributions that package it.
\emph{Dissemination} occurs through official binary packages, distribution
repositories, and container base images. The \emph{user base} spans
essentially every sector that runs open-source infrastructure, including the
health-care, financial, utility, and government deployments that this
proposal targets. The \emph{contributor base}, in contrast, is comparatively
small, and it consists of a core maintainer group, a set of occasional patch
contributors, and a large third-party module ecosystem that is maintained
outside the project. This asymmetry between an enormous user base and a concentrated maintainer base is precisely the structural risk PESOSE was created to address.

\paragraph{ZMap: discovering what is actually running.}
ZMap is a single-packet, stateless network scanner that has been released
under the Apache~2.0 license and that is maintained by an academic and
industrial consortium~\cite{zmap,zmapsrc}. Its stateless design allows a
single commodity machine to sweep enormous address spaces far faster than
connection-oriented scanners, and the companion ZGrab tool performs
application-layer handshakes that identify service software and
versions~\cite{zgrab}. The \emph{development and testing model} is a
pull-request-and-review process conducted in the open on GitHub, with unit
and integration tests that are run on every pull request and with tagged
releases that are packaged for major platforms. \emph{Dissemination} occurs
through the source repository, distribution packages, and container images,
and ZMap underpins commercial and non-profit internet-measurement services as
well as a substantial body of peer-reviewed research. The \emph{user base} includes network operators who perform asset inventory,
incident responders, attack-surface-management services, and academic
measurement groups, and the \emph{contributor base} is an active but modest
community centered on the maintaining institutions.

ZMap was designed for internet-wide measurement from outside a network, where
the marginal effect of a probe on any single host is negligible. This project
requires the opposite, that is, repeated discovery \emph{inside} a production
network whose hosts may be fragile medical, industrial, or embedded devices.
Applied naively, the very property that makes ZMap valuable, its speed, is what makes it unsafe here. Section~\ref{sec:thrust1} describes the extensions we will contribute upstream to make ZMap a first-class instrument for safe, budgeted, and information-rich discovery on production infrastructure.

\paragraph{The ecosystem view.}
These two products bracket the problem that \sysname addresses, as ZMap
determines what an operator can \emph{know} about an infrastructure without
disturbing it, and nginx determines what an attacker can \emph{reach} once
inside it. Both sit within a broader ecosystem that \sysname must also
capture, replicate, and analyze, including PostgreSQL~\cite{postgresql}, Keycloak~\cite{keycloak},
RabbitMQ~\cite{rabbitmq}, and Istio~\cite{istio}. Our working relationships with the maintainers and users of nginx and ZMap, documented in the accompanying letters of collaboration, give the project a direct path from a finding produced on a twin to a fix that reaches every downstream user.
